% =====================================================================
\section{Rendering Pipeline}\label{sec:pipeline}
The frame is rendered in linear light into a multisampled HDR target and post-processed. The passes in
Fig.~\ref{fig:frame} correspond one-to-one to the GPU timer labels in the code (\code{timer\_.begin(...)} in
\code{Renderer.cpp}/\code{RendererFrame.cpp}): \emph{uav camera (CV), shadow maps, pre-pass, ssao, volumetric fog, scene,
combine + ssr, taa, exposure / agc, lens flare, depth of field, resolve, bloom, composite}. The pipeline is the v3 pipeline;
v4 adds the effects renderer, the indoor mask, the 16-light ranking with dynamic lights and the seven surfaces.

\begin{figure}[h]
\centering
\begin{tikzpicture}[node distance=2.6mm and 3.2mm, every node/.style={font=\scriptsize},
  box/.style={draw,rounded corners=2pt,minimum height=7mm,align=center,fill=orange!6},
  arr/.style={-{Latex[length=1.5mm]},thick}]
  \node[box] (s) {shadow maps\\moon $\times2$, spot};
  \node[box,right=of s] (p) {pre-pass\\depth+normal};
  \node[box,right=of p] (a) {SSAO};
  \node[box,right=of a] (f) {volumetric\\fog};
  \node[box,right=of f] (m) {HDR scene\\+ effects};
  \node[box,right=of m] (c) {combine\\fog + SSR};
  \node[box,below=of s] (t) {TAA};
  \node[box,right=of t] (e) {exposure /\\thermal AGC};
  \node[box,right=of e] (b) {bloom};
  \node[box,right=of b] (l) {lens flare\\+ streak};
  \node[box,right=of l] (d) {depth of\\field};
  \node[box,right=of d] (o) {composite\\ACES / FLIR};
  \foreach \x/\y in {s/p,p/a,a/f,f/m,m/c,t/e,e/b,b/l,l/d,d/o} \draw[arr] (\x)--(\y);
  \draw[arr] (c.south) -- ++(0,-2mm) -| (t.north);
\end{tikzpicture}
\caption{The frame. A separate raw heat render from the gimbal camera ($320\times240$, no post-processing) feeds the
thermal tracker.}
\label{fig:frame}
\end{figure}

\subsection{Shadows: cascades, static cache and dynamic flags}
The moon uses two orthographic cascades in a depth-texture array. The \emph{near cascade} is fitted around the
bounding sphere of the view-frustum slice $[z_n,\text{split}]$ (split 10\,m LOW, 12\,m HIGH, 14\,m ULTRA/EPIC) and is redrawn
every frame; the \emph{world cascade} is fitted around a sphere of radius $r=\texttt{kBaseRadius}+3=39$\,m centred at
$(0,1.5,0)$, i.e.\ about $78\,\text{m}/4096=1.9$\,cm per texel on EPIC and 3.8\,cm at 2048. The cascade centre is snapped to whole texels in light space,
$c'=\lfloor c/\Delta\rfloor\Delta$, $\Delta=2r/\text{res}$, so edges do not shimmer (T34). The world cascade never
moves, so its \textbf{static} casters are rendered once into a cache that is copied into the array each frame; only
\textbf{dynamic} nodes are drawn on top. A node is dynamic if \code{SceneNode::dynamic} is set (the Scene marks the UAV,
the rovers and every controllable object other than the drone controller's own root) or it belongs to the UAV or
rover. The cache key is the explode value plus the foliage setting, so the static layer is rebuilt only if the geometry
set changes. The searchlight has a perspective shadow map along the lamp axis (it excludes the UAV, the lamp is inside
it). The backdrop, blur discs, the beam and the moon do not cast shadows. Lookups use a normal-offset bias and a rotated
Poisson disc of 8 (HIGH) or 16 (ULTRA/EPIC) hardware-PCF taps.

\subsection{SSAO, volumetric fog, bloom}
\textbf{SSAO} samples a normal-oriented hemisphere (8 samples HIGH, 16 EPIC) around each pixel of the depth/normal
pre-pass (half resolution, full resolution on EPIC) and is blurred depth-aware. \textbf{Volumetric fog} ray-marches a
height fog $\sigma(y)=\sigma_0e^{-0.22y}$ lit by the moon, the beacon and the practical lights with single
scattering, $L\mathrel{+}=T\,L_\text{in}(1-e^{-\sigma\Delta s})$, $T\mathrel{*}=e^{-\sigma\Delta s}$ (24 steps HIGH, 48
ULTRA, 64 EPIC at half resolution). The searchlight shaft is a thin cone that a coarse march would miss, so the ray is
intersected analytically with the cone $(\vect d\cdot\vect a)^2=\cos^2\gamma\,|\vect x-\vect s|^2$ and only that
interval is sampled, using the spot shadow map. Underground the fog is multiplied by $1-0.85\,\text{mask}$.
\textbf{Bloom} uses a soft-knee threshold, a 13-tap downsample chain and a tent upsample.

\subsection{TAA, SSR, DOF, motion blur, flare, auto-exposure, tone mapping}
\begin{itemize}
  \item \textbf{TAA}: Halton jitter, reprojection from the full-resolution depth, Catmull--Rom history, YCoCg neighbourhood variance
  clipping, cut detection; captures accumulate 8 frames. Switch: \texttt{F6}.
  \item \textbf{SSR}: screen-space reflections for puddles and wet ground from the surface buffer (second render target
  with normal and reflectivity), 48 steps (64 on EPIC). Switch: \texttt{F7}.
  \item \textbf{Lens flare}: ghosts, halo and an anamorphic streak from the bloom chain; off indoors. Switch: \texttt{F8}.
  \item \textbf{Auto-exposure}: log-average luminance, temporally adapted, clamped to $\pm1.5$\,EV around a key of $\log_2 L=-5.2$ so that
  a night stays dark. Switch: \texttt{F9}; \texttt{+}/\texttt{-} adjust manually.
  \item \textbf{Depth of field}: golden-angle bokeh gather; \textbf{motion blur}: from depth reprojection (cinematic shots and video).
  \texttt{F10} toggles DOF.
  \item \textbf{ACES}: fitted curve $\frac{x(2.51x+0.03)}{x(2.43x+0.59)+0.14}$, sRGB, scotopic (Purkinje) shift and split toning,
  vignette and grain.
\end{itemize}

\subsection{Thermal sensor model}
In THERMAL view (key \texttt{V}; \texttt{Y} switches between the first-person gimbal view and the external view) heat is
rendered as a scalar and the sensor pass applies: automatic gain control from the frame's heat range (temporally
smoothed), digital detail enhancement against a wide local mean, a palette (\texttt{T}: ironbow, white-hot, rainbow,
black-hot), temporal noise, fixed-pattern column/row offsets, a few dead pixels, a narcissus vignette, heat shimmer above
hot parts and a periodic NUC shutter freeze (\texttt{Q} forces one). The zoom (mouse wheel) is WFOV/MFOV/NFOV at
40/20/8$^\circ$. The FLIR-style HUD shows compass, gimbal scales, telemetry, target box with slant range and track state.
In v4 the explosion particles and the scorch marks write a heat value too (below), so a fireball is a bright blob in the
thermal view and a wreck stays hot while it cools.

\subsection{Effects renderer}
Explosions and weapons are drawn by \code{EffectsRenderer} (\code{rendering/Effects.cpp}). The sprites are instanced
quads taken from the lab's unit cube (only its $+Z$ face is used; the other faces are collapsed in the vertex shader), one
instance being $4\times$\code{vec4}. All kinds share one program and one premultiplied-alpha blend state
\code{glBlendFunc(ONE, ONE\_MINUS\_SRC\_ALPHA)}: smoke writes (lit colour$\cdot\alpha$, $\alpha$), emissive sprites write
(HDR colour, $\sim0$), so ``additive'' and ``alpha'' differ only in the alpha value. The draw order is smoke, trail and
dust sorted back to front, then the emissive group (fire, flash, sparks, embers, tracers, shock waves), again sorted, because
the fire's soot edge needs the order. Scorch decals and projectile bodies (missile, shell) are drawn after the solids with the
depth test. The soft-particle depth fade exists in the shader (\code{uHasSceneDepth}) but the renderer currently sets it to 0.
The eight particle types are listed in Table~\ref{tab:fx}; the particle cap is 4\,000 (T104), and low-priority smoke, trail and
dust are dropped first.

\begin{table}[h]
\centering\small
\caption{Effect particle types (\code{EffectType}, \code{EffectsData.h}) and their role. Brightness is scaled at spawn in
\code{Combat::spawn} (flash $\times0.12$, fire $\times0.14$, spark and ember $\times0.45$) to land at a few HDR units under the
night exposure.}
\label{tab:fx}
\begin{tabular}{@{}rllp{9.3cm}@{}}
\toprule
Id & Type & Blend & Role\\\midrule
0 & Fire & emissive & fireball puff, orange $\to$ red $\to$ dark, soot edge\\
1 & Smoke & lit, alpha & grows, dark grey (explosion) or light grey (exhaust)\\
2 & Spark & emissive & tiny hot streak along its velocity\\
3 & Flash & emissive & very short bright sphere (detonation, muzzle flash)\\
4 & Shockwave & emissive & expanding faint ring on the ground\\
5 & Trail & lit, alpha & thin missile and shell exhaust puff\\
6 & Dust & lit, alpha & brown dust kicked by blasts, tracks, rotor wash\\
7 & Ember & emissive & glowing debris speck that falls and cools\\\bottomrule
\end{tabular}
\end{table}

\subsection{Quality presets}
\begin{table}[h]
\centering\small
\caption{Presets from \code{qualityProfile()} (RendererFrame.cpp); EPIC is auto-selected on NVIDIA / discrete GPUs
(override with \cli{quality}). Values not listed keep the HIGH defaults.}
\label{tab:presets}
\begin{tabular}{@{}lrrrrrrrr@{}}
\toprule
Preset & MSAA & Moon / spot map & SSAO & PCF & Fog steps & Grass blades & Rain drops & Fireflies\\\midrule
LOW & 2 & 1024 / off & off & 8 & off & 14\,000 & 6\,000 & 60\\
HIGH & 4 & 2048 / 1024 & 8 & 8 & 24 & 35\,000 & 14\,000 & 200\\
ULTRA & 8 & 4096 / 2048 & 16 & 16 & 48 & 140\,000 & 36\,000 & 400\\
EPIC & 4 + TAA & 4096 / 2048 & 16, full res & 16 & 64 & 450\,000 & 60\,000 & 700\\\bottomrule
\end{tabular}
\end{table}
Also: SSR, DOF and motion blur are on for ULTRA and EPIC; TAA on for HIGH and above; foliage density 0.2 / 0.5 / 0.7 / 1.0.
